\begin{abstract}
针对六区域异构 AI 任务与新能源、储能、电网的耦合调度，本文构建“闭环需求预测—有限任务邻域—离散 SOC 动态规划—按情景交替搜索”的可复算框架。任务侧用连续区间重叠区分瞬时容量与小时电量，能源侧统一约束购售电、充放电、能量守恒及终端 SOC。问题一在独立验证集上从季节朴素与 Huber 候选中选模，18 条序列中 17 条选择 Huber，平均 MAE、RMSE 和 WAPE 分别为 25.5892 GPU、34.5070 GPU 和 0.8231。问题二每个方向评价 10000 个候选；平衡方案成本为 $-4.218017\times10^8$ 元，碳排放 16.5160 tCO$_2$，迁移 561 个任务。问题三的 7 个策略形成 6 条唯一轨迹，平衡策略将成本由 $-4.198203\times10^8$ 元降至 $-4.585846\times10^8$ 元，碳排放由 8.8506 降至 0.2241 tCO$_2$。问题四对 17 个情景动态重建任务邻域；联合方案的成本、碳排放和峰值购电分别为 $-4.587163\times10^8$ 元、0.2684 tCO$_2$ 和 0.3644 MW。双运行除自然耗时字段外完全一致。全部结论均为有限搜索 best-found 的场景可行结果，不构成全局最优或完整 Pareto 前沿。
\end{abstract}
\keywords{异构 AI 任务；闭环预测；算电协同；储能动态规划；多目标调度}
